\documentclass[12pt,a4paper]{report}
\usepackage[bahasa]{babel}
\usepackage[utf8]{inputenc}
\usepackage[T1]{fontenc}
\usepackage{geometry}
\geometry{a4paper, margin=2.5cm, top=3cm, bottom=2.5cm, headheight=14.5pt}
\usepackage{xcolor}
\usepackage[most]{tcolorbox}
\usepackage{graphicx}
\usepackage{booktabs}
\usepackage{tabularx}
\usepackage{enumitem}
\usepackage[expansion=false]{microtype}
\usepackage{titlesec}
\usepackage{fancyhdr}
\usepackage{amsmath}
\usepackage{amssymb}

% Warna Korporat PLN
\definecolor{plnblue}{RGB}{0,75,135}
\definecolor{plnorange}{RGB}{255,107,53}
\definecolor{fillgreen}{RGB}{240,255,240}

\usepackage{hyperref}
\hypersetup{
    colorlinks=true,
    linkcolor=plnblue,
    urlcolor=plnblue,
    citecolor=plnblue
}

% Header & Footer
\pagestyle{fancy}
\fancyhf{}
\fancyhead[L]{\small\textbf{Pre-Test Modul 4.3}}
\fancyhead[R]{\small Level 4 (Mastery) -- PLN}
\fancyfoot[C]{\small\thepage}
\renewcommand{\headrulewidth}{0.4pt}

% Format Heading
\titleformat{\chapter}[display]{\normalfont\huge\bfseries\color{plnblue}}{\chaptertitlename\ \thechapter}{20pt}{\Huge}
\titleformat{\section}{\Large\bfseries\color{plnblue}}{}{0em}{}
\titleformat{\subsection}{\large\bfseries\color{plnblue}}{}{0em}{}
\titlespacing*{\chapter}{0pt}{-30pt}{20pt}
\titlespacing*{\section}{0pt}{12pt}{8pt}
\titlespacing*{\subsection}{0pt}{10pt}{6pt}

% Custom Boxes
\newtcolorbox{plnbox}[1][]{%
  colback=plnblue!5!white, colframe=plnblue, fonttitle=\bfseries,
  title=#1, sharp corners, boxrule=0.8pt, breakable
}
\newtcolorbox{warnbox}{%
  colback=plnorange!10, colframe=plnorange, fonttitle=\bfseries,
  title=Petunjuk Pengerjaan, sharp corners, breakable
}
\newtcolorbox{answerbox}{%
  colback=fillgreen, colframe=green!50!black, fonttitle=\bfseries,
  title=Kunci Jawaban (Khusus Fasilitator), sharp corners, breakable
}

% List styling
\setlist[itemize]{leftmargin=*, topsep=2pt, itemsep=2pt}
\setlist[enumerate]{leftmargin=*, topsep=2pt, itemsep=2pt}

% Helper: answer lines
\newcommand{\ansline}[1][5.5cm]{\underline{\hspace{#1}}}

\begin{document}

% ================= COVER =================
\begin{titlepage}
\centering
\vspace*{1.5cm}
{\Huge\bfseries\color{plnblue} PRE-TEST}\\[0.4cm]
{\LARGE\bfseries Modul 4.3: Optimalisasi Prosedur dan}\\[0.2cm]
{\LARGE\bfseries Alur Kerja Terintegrasi}\\[0.8cm]
\rule{0.8\linewidth}{0.6pt}\\[0.6cm]
{\large Program Penguatan Kompetensi Tata Kelola Kearsipan}\\
{\large Level 4 (Mastery)}\\[1.2cm]
{\large Disusun Khusus untuk: \textbf{PT PLN (Persero)}}\\
{\large Target Peserta: Manajemen Atas (MA)}\\[2cm]

\begin{plnbox}[Tujuan Pre-Test]
\centering
Pre-test ini \textbf{bukan} penilaian kelulusan. Tujuannya adalah \textbf{mengukur pemahaman awal} peserta agar fasilitator dapat menyesuaikan kedalaman materi dan pendampingan selama sesi pelatihan. Jawablah sejujurnya berdasarkan pemahaman Anda saat ini.
\end{plnbox}

\vspace{1cm}
\begin{figure}[h]
\centering
\includegraphics[width=4cm]{example-image}
\end{figure}
\vfill
{\small Waktu Pengerjaan: 15 Menit \quad|\quad Sifat: Diagnostik (Bukan Penilaian Kelulusan)}
\end{titlepage}

% ================= IDENTITAS PESERTA =================
\newpage
\thispagestyle{fancy}

\begin{center}
{\LARGE\bfseries\color{plnblue} LEMBAR JAWABAN PRE-TEST}\\[0.3cm]
{\large Modul 4.3 -- Level 4 (Mastery)}
\end{center}

\vspace{0.6cm}

\renewcommand{\arraystretch}{1.8}
\noindent
\begin{tabularx}{\textwidth}{@{} X X @{}}
\textbf{Nama Peserta:} \ansline & \textbf{Unit Kerja:} \ansline \\
\textbf{Jabatan:} \ansline & \textbf{Tanggal:} \ansline \\
\end{tabularx}
\renewcommand{\arraystretch}{1.0}

\vspace{0.4cm}

\begin{warnbox}
\begin{itemize}[nosep]
    \item Waktu pengerjaan: \textbf{15 menit}. Kerjakan seluruh bagian (A--C).
    \item \textbf{Tidak ada jawaban salah yang dihukum.} Jika ragu, tetap jawab sesuai perkiraan terbaik Anda.
    \item Untuk Bagian A, lingkari atau beri tanda \textbf{X} pada jawaban yang dipilih.
    \item Untuk Bagian B, beri tanda centang ($\checkmark$) pada kolom yang sesuai.
    \item Untuk Bagian C, tulis jawaban singkat pada ruang yang disediakan.
    \item Hasil pre-test akan dibandingkan dengan post-test untuk mengukur efektivitas pembelajaran.
\end{itemize}
\end{warnbox}

\vspace{0.3cm}

\renewcommand{\arraystretch}{1.3}
\begin{center}
\small
\begin{tabularx}{0.85\textwidth}{@{} X c c c @{}}
\toprule
\textbf{Bagian} & \textbf{Jumlah Soal} & \textbf{Bobot/Soal} & \textbf{Total Poin} \\
\midrule
A. Pilihan Ganda (Pengetahuan Dasar) & 10 soal & 4 poin & 40 poin \\
B. Penilaian Diri (\textit{Self-Assessment}) & 5 pernyataan & 2 poin & 10 poin \\
C. Refleksi Pengalaman Kerja & 5 soal & 10 poin & 50 poin \\
\midrule
\textbf{TOTAL} & \textbf{20 soal} & & \textbf{100 poin} \\
\bottomrule
\end{tabularx}
\end{center}
\renewcommand{\arraystretch}{1.0}

% ================= BAGIAN A: PILIHAN GANDA =================
\newpage
\section{BAGIAN A: Pilihan Ganda -- Pengetahuan Dasar (10 soal $\times$ 4 poin = 40 poin)}
\textit{Soal-soal berikut menguji pemahaman prasyarat Anda tentang tata kelola kearsipan dan proses bisnis. Lingkari \textbf{SATU} jawaban yang paling tepat.}

\vspace{0.3cm}

\begin{enumerate}[leftmargin=*, itemsep=0.5cm]

% --- Soal 1: Regulasi Dasar ---
\item Regulasi utama yang mengatur penyelenggaraan sistem dan transaksi elektronik di Indonesia, termasuk keabsahan tanda tangan elektronik, adalah:
\begin{enumerate}[label=\Alph*., nosep]
    \item UU No. 14 Tahun 2008 tentang Keterbukaan Informasi Publik
    \item PP No. 71 Tahun 2019 tentang Penyelenggaraan Sistem dan Transaksi Elektronik (PSTE)
    \item Perka ANRI No. 28 Tahun 2012 tentang Pedoman Tata Naskah Dinas
    \item UU No. 25 Tahun 2009 tentang Pelayanan Publik
\end{enumerate}

% --- Soal 2: Definisi Arsip ---
\item Menurut UU 43/2009 tentang Kearsipan, yang dimaksud ``arsip'' adalah:
\begin{enumerate}[label=\Alph*., nosep]
    \item Hanya dokumen kertas yang disimpan di lemari/brankas
    \item Rekaman kegiatan atau peristiwa dalam berbagai bentuk dan media, yang dibuat dan diterima oleh lembaga
    \item Salinan digital dari dokumen yang sudah dihancurkan
    \item File komputer yang tersimpan di server pusat data
\end{enumerate}

% --- Soal 3: Konsep Proses Bisnis ---
\item Istilah ``\textit{bottleneck}'' dalam konteks alur kerja (proses bisnis) merujuk pada:
\begin{enumerate}[label=\Alph*., nosep]
    \item Titik di mana proses mengalami hambatan/penumpukan sehingga memperlambat keseluruhan alur
    \item Jumlah maksimum dokumen yang dapat disimpan di server
    \item Tahap akhir dari suatu proses persetujuan
    \item Teknik untuk mempercepat pemrosesan data dengan kompresi file
\end{enumerate}

% --- Soal 4: Metadata ---
\item Metadata arsip adalah:
\begin{enumerate}[label=\Alph*., nosep]
    \item Isi/konten utama dari dokumen itu sendiri
    \item Informasi deskriptif tentang arsip (seperti judul, tanggal, pencipta, klasifikasi) yang memudahkan pencarian dan pengelolaan
    \item Salinan cadangan (\textit{backup}) dari arsip asli
    \item Tanda tangan digital yang melekat pada dokumen
\end{enumerate}

% --- Soal 5: Retensi Arsip ---
\item Jadwal Retensi Arsip (JRA) menentukan:
\begin{enumerate}[label=\Alph*., nosep]
    \item Berapa banyak salinan arsip yang harus dibuat
    \item Siapa saja yang boleh mengakses arsip tertentu
    \item Berapa lama arsip harus disimpan sebelum dipindahkan, dimusnahkan, atau diserahkan ke lembaga kearsipan
    \item Format file yang wajib digunakan untuk menyimpan arsip digital
\end{enumerate}

% --- Soal 6: SOP ---
\item Fungsi utama Standar Operasional Prosedur (SOP) dalam pengelolaan arsip adalah:
\begin{enumerate}[label=\Alph*., nosep]
    \item Menggantikan seluruh peran manusia dengan otomasi komputer
    \item Menjadi panduan baku yang menstandarkan langkah kerja, peran, dan tanggung jawab agar proses berjalan konsisten dan teraudit
    \item Mengurangi jumlah pegawai yang dibutuhkan pada unit kearsipan
    \item Menyimpan password dan kredensial akses sistem
\end{enumerate}

% --- Soal 7: Integrasi Sistem ---
\item Dalam konteks organisasi, ``integrasi sistem'' secara umum berarti:
\begin{enumerate}[label=\Alph*., nosep]
    \item Mengganti seluruh aplikasi lama dengan satu aplikasi baru tunggal
    \item Menghubungkan beberapa sistem/aplikasi yang berbeda agar dapat bertukar data dan bekerja secara terkoordinasi
    \item Memindahkan seluruh data ke \textit{cloud computing}
    \item Membeli lisensi perangkat lunak dari satu vendor yang sama
\end{enumerate}

% --- Soal 8: Tanda Tangan Elektronik ---
\item Tanda tangan elektronik yang bersertifikat menurut hukum Indonesia:
\begin{enumerate}[label=\Alph*., nosep]
    \item Tidak memiliki kekuatan hukum sama sekali
    \item Hanya berlaku untuk dokumen internal perusahaan, tidak untuk eksternal
    \item Memiliki kekuatan hukum dan akibat hukum yang sah, setara dengan tanda tangan basah
    \item Hanya berlaku jika disertai materai elektronik
\end{enumerate}

% --- Soal 9: Siklus Hidup Arsip ---
\item Urutan siklus hidup arsip (\textit{records lifecycle}) yang benar adalah:
\begin{enumerate}[label=\Alph*., nosep]
    \item Penyimpanan $\rightarrow$ Penciptaan $\rightarrow$ Penggunaan $\rightarrow$ Pemusnahan
    \item Penciptaan $\rightarrow$ Penggunaan/Pemeliharaan $\rightarrow$ Penyusutan (Pemusnahan/Penyerahan)
    \item Pemusnahan $\rightarrow$ Penciptaan $\rightarrow$ Penyimpanan $\rightarrow$ Distribusi
    \item Pengadaan $\rightarrow$ Instalasi $\rightarrow$ Operasi $\rightarrow$ Dekomisi
\end{enumerate}

% --- Soal 10: Audit Kearsipan ---
\item Salah satu risiko utama jika arsip tidak dikelola dengan baik di lingkungan BUMN seperti PLN adalah:
\begin{enumerate}[label=\Alph*., nosep]
    \item Kecepatan internet menurun
    \item Temuan audit oleh BPK/auditor eksternal terkait ketidaklengkapan jejak dokumen
    \item Tagihan listrik pelanggan mengalami kenaikan
    \item Kapasitas pembangkit listrik berkurang
\end{enumerate}

\end{enumerate}

% ================= BAGIAN B: PENILAIAN DIRI =================
\newpage
\section{BAGIAN B: Penilaian Diri / \textit{Self-Assessment} (5 pernyataan $\times$ 2 poin = 10 poin)}
\textit{Beri tanda centang ($\checkmark$) pada kolom yang paling menggambarkan tingkat pemahaman Anda \textbf{saat ini} terhadap topik berikut. Jawablah dengan jujur -- tidak ada jawaban benar atau salah.}

\vspace{0.4cm}

\renewcommand{\arraystretch}{1.6}
\begin{center}
\small
\begin{tabularx}{\textwidth}{@{} c >{\raggedright\arraybackslash}X c c c c @{}}
\toprule
\textbf{No} & \textbf{Pernyataan Topik} & \rotatebox{90}{\textbf{Belum Tahu}} & \rotatebox{90}{\textbf{Pernah Dengar}} & \rotatebox{90}{\textbf{Cukup Paham}} & \rotatebox{90}{\textbf{Sangat Paham}} \\
\midrule
1 & Saya memahami perbedaan antara evaluasi prosedur kearsipan dan audit infrastruktur IT. & $\square$ & $\square$ & $\square$ & $\square$ \\[0.3cm]
2 & Saya mengetahui apa itu \textit{Value Stream Mapping} dan kegunaannya untuk mengidentifikasi pemborosan dalam alur kerja. & $\square$ & $\square$ & $\square$ & $\square$ \\[0.3cm]
3 & Saya pernah melihat atau menggunakan diagram BPMN (\textit{Business Process Model and Notation}) untuk menggambarkan alur kerja. & $\square$ & $\square$ & $\square$ & $\square$ \\[0.3cm]
4 & Saya mengetahui apa yang dimaksud dengan ``pemicu otomatis'' dalam konteks integrasi antar-sistem (misalnya: satu status di Sistem A memicu aksi di Sistem B). & $\square$ & $\square$ & $\square$ & $\square$ \\[0.3cm]
5 & Saya dapat menjelaskan mengapa metadata arsip penting untuk temu kembali (\textit{retrieval}) dan kepatuhan audit. & $\square$ & $\square$ & $\square$ & $\square$ \\
\bottomrule
\end{tabularx}
\end{center}
\renewcommand{\arraystretch}{1.0}

\vspace{0.4cm}

\begin{plnbox}[Catatan]
\small
Hasil penilaian diri ini akan digunakan oleh fasilitator untuk:
\begin{itemize}[nosep]
    \item Mengidentifikasi topik yang memerlukan penjelasan lebih mendalam di kelas
    \item Membentuk kelompok workshop yang heterogen (campuran peserta berpengalaman dan pemula)
    \item Menyesuaikan kecepatan penyampaian materi per bab
\end{itemize}
\end{plnbox}


% ================= BAGIAN C: REFLEKSI PENGALAMAN =================
\newpage
\section{BAGIAN C: Refleksi Pengalaman Kerja (5 soal $\times$ 10 poin = 50 poin)}
\textit{Jawab berdasarkan pengalaman nyata di unit kerja Anda. Jika belum memiliki pengalaman terkait, tuliskan ``Belum pernah mengalami'' dan jelaskan apa yang Anda perkirakan akan terjadi. Jawab maksimal 3--4 kalimat per soal.}

\vspace{0.3cm}

\begin{enumerate}[leftmargin=*, itemsep=0.3cm]

% --- C1: Identifikasi Masalah ---
\item Sebutkan \textbf{1 masalah utama} yang sering Anda temui dalam proses pengelolaan arsip/dokumen di unit kerja Anda saat ini. Apa dampaknya terhadap operasional sehari-hari?

\vspace{0.2cm}
\noindent\rule{\textwidth}{0.3pt}
\vspace{0.5cm}
\noindent\rule{\textwidth}{0.3pt}
\vspace{0.5cm}
\noindent\rule{\textwidth}{0.3pt}
\vspace{0.5cm}
\noindent\rule{\textwidth}{0.3pt}

% --- C2: Pengalaman Duplikasi ---
\item Pernahkah Anda atau staf Anda harus \textbf{menginput data yang sama ke lebih dari satu sistem/aplikasi} untuk satu proses yang sama? Jika ya, ceritakan situasinya secara singkat. Jika tidak, menurut Anda mengapa redundansi input data bisa merugikan?

\vspace{0.2cm}
\noindent\rule{\textwidth}{0.3pt}
\vspace{0.5cm}
\noindent\rule{\textwidth}{0.3pt}
\vspace{0.5cm}
\noindent\rule{\textwidth}{0.3pt}
\vspace{0.5cm}
\noindent\rule{\textwidth}{0.3pt}

% --- C3: Tantangan Pencarian ---
\item Ketika ada permintaan mendadak untuk menemukan arsip tertentu (misalnya saat audit atau permintaan data dari manajemen), berapa lama waktu yang biasanya dibutuhkan? Apa kesulitan utamanya?

\vspace{0.2cm}
\noindent\rule{\textwidth}{0.3pt}
\vspace{0.5cm}
\noindent\rule{\textwidth}{0.3pt}
\vspace{0.5cm}
\noindent\rule{\textwidth}{0.3pt}
\vspace{0.5cm}
\noindent\rule{\textwidth}{0.3pt}

% --- C4: Pengalaman Persetujuan ---
\item Bagaimana proses persetujuan/validasi dokumen di unit kerja Anda? Apakah sudah menggunakan persetujuan digital, atau masih memerlukan tanda tangan fisik/basah? Apa kendala yang paling sering muncul?

\vspace{0.2cm}
\noindent\rule{\textwidth}{0.3pt}
\vspace{0.5cm}
\noindent\rule{\textwidth}{0.3pt}
\vspace{0.5cm}
\noindent\rule{\textwidth}{0.3pt}
\vspace{0.5cm}
\noindent\rule{\textwidth}{0.3pt}

% --- C5: Harapan Pembelajaran ---
\item Apa \textbf{1 hal utama} yang ingin Anda pelajari atau peroleh dari pelatihan modul ini? Masalah spesifik apa di unit kerja Anda yang Anda harapkan bisa diatasi setelah mengikuti pelatihan?

\vspace{0.2cm}
\noindent\rule{\textwidth}{0.3pt}
\vspace{0.5cm}
\noindent\rule{\textwidth}{0.3pt}
\vspace{0.5cm}
\noindent\rule{\textwidth}{0.3pt}
\vspace{0.5cm}
\noindent\rule{\textwidth}{0.3pt}

\end{enumerate}

% ================= KUNCI JAWABAN (FASILITATOR) =================
\newpage
\begin{center}
{\LARGE\bfseries\color{red!70!black} KUNCI JAWABAN \& PANDUAN PENILAIAN}\\[0.2cm]
{\large\color{red!70!black}\textit{(Dokumen Rahasia -- Khusus Fasilitator)}}
\end{center}

\vspace{0.5cm}

\begin{answerbox}
\textbf{PENTING:} Halaman ini \textbf{WAJIB DIPISAHKAN} dari lembar soal peserta. Pre-test bersifat diagnostik -- fasilitator menggunakan hasil ini untuk \textbf{memetakan keragaman pengetahuan awal kelas}, bukan untuk menentukan kelulusan. Bandingkan hasil pre-test dengan post-test di akhir sesi untuk mengukur \textit{learning gain}.
\end{answerbox}

\vspace{0.5cm}

% --- Kunci Bagian A ---
\subsection{Kunci Jawaban Bagian A: Pilihan Ganda (40 poin)}

\renewcommand{\arraystretch}{1.4}
\begin{center}
\begin{tabularx}{0.95\textwidth}{@{} c c X @{}}
\toprule
\textbf{No} & \textbf{Jawaban} & \textbf{Konsep Prasyarat yang Diuji} \\
\midrule
1 & \textbf{B} & Pengetahuan regulasi PSTE -- fondasi hukum tanda tangan elektronik \\
2 & \textbf{B} & Definisi arsip menurut UU 43/2009 -- cakupan tidak terbatas pada kertas \\
3 & \textbf{A} & Terminologi proses bisnis dasar -- hambatan proses alur kerja \\
4 & \textbf{B} & Konsep metadata -- informasi deskriptif tentang arsip \\
5 & \textbf{C} & Jadwal Retensi Arsip -- durasi simpan dan nasib akhir arsip \\
6 & \textbf{B} & Fungsi SOP -- standarisasi langkah, peran, dan konsistensi \\
7 & \textbf{B} & Integrasi sistem -- koneksi antar-aplikasi untuk pertukaran data \\
8 & \textbf{C} & Keabsahan TTD elektronik bersertifikat -- setara TTD basah \\
9 & \textbf{B} & Siklus hidup arsip -- penciptaan, penggunaan, penyusutan \\
10 & \textbf{B} & Risiko tata kelola arsip buruk di BUMN -- temuan audit BPK \\
\bottomrule
\end{tabularx}
\end{center}
\renewcommand{\arraystretch}{1.0}

\vspace{0.5cm}

% --- Panduan Bagian B ---
\subsection{Panduan Penilaian Bagian B: Penilaian Diri (10 poin)}

\begin{plnbox}[Cara Skoring \& Interpretasi]
\small
Penilaian diri \textbf{tidak dinilai benar/salah}. Skor diberikan berdasarkan \textbf{kelengkapan pengisian} (2 poin per baris yang diisi). Fasilitator menggunakan hasil ini sebagai peta diagnostik:

\vspace{0.3cm}
\renewcommand{\arraystretch}{1.3}
\begin{tabularx}{\textwidth}{@{} c X @{}}
\toprule
\textbf{Kolom Dominan} & \textbf{Implikasi bagi Fasilitator} \\
\midrule
Mayoritas ``Belum Tahu'' & Peserta memerlukan pengantar konseptual lebih panjang di Bab 1--2. Berikan contoh konkret PLN lebih banyak. \\
Mayoritas ``Pernah Dengar'' & Peserta memiliki \textit{awareness} namun belum \textit{hands-on}. Fokuskan pada aktivitas workshop dan studi kasus. \\
Mayoritas ``Cukup/Sangat Paham'' & Percepat pengantar, alokasikan waktu lebih banyak untuk diskusi lanjutan dan desain BPMN/SOP. \\
\bottomrule
\end{tabularx}
\renewcommand{\arraystretch}{1.0}
\end{plnbox}

\subsection{Panduan Penilaian Bagian C: Refleksi Pengalaman (50 poin)}

\textit{Bagian ini dinilai berdasarkan \textbf{kedalaman refleksi}, bukan ketepatan terminologi. Peserta yang belum berpengalaman tetap mendapat kredit jika menunjukkan kemampuan berpikir analitis.}

\vspace{0.3cm}

\renewcommand{\arraystretch}{1.4}
\begin{center}
\begin{tabularx}{0.95\textwidth}{@{} c X X @{}}
\toprule
\textbf{No} & \textbf{Komponen Jawaban Berkualitas} & \textbf{Distribusi Poin} \\
\midrule
C1 & Menyebutkan masalah spesifik (bukan generik) + menjelaskan dampak operasional nyata & 5 poin masalah + 5 poin dampak \\
\hline
C2 & Mendeskripsikan situasi duplikasi konkret atau menjelaskan logis mengapa redundansi merugikan (waktu, inkonsistensi, risiko error) & 5 poin deskripsi + 5 poin analisis rugi \\
\hline
C3 & Menyebutkan estimasi waktu pencarian + mengidentifikasi penyebab kesulitan (tidak terindeks, tersebar, metadata buruk) & 5 poin waktu + 5 poin penyebab \\
\hline
C4 & Menggambarkan proses persetujuan aktual di unit + mengidentifikasi kendala (TTD basah, atasan di luar kantor, tidak ada SLA) & 5 poin deskripsi + 5 poin kendala \\
\hline
C5 & Menyebutkan harapan spesifik dan realistis (bukan generik seperti ``ingin belajar'') + menghubungkan dengan masalah nyata di unit kerja & 5 poin harapan + 5 poin koneksi masalah \\
\bottomrule
\end{tabularx}
\end{center}
\renewcommand{\arraystretch}{1.0}

\vspace{0.5cm}

% --- Panduan Perbandingan Pre-Post ---
\subsection{Panduan Perbandingan Pre-Test vs Post-Test}

\begin{plnbox}[Formula \textit{Learning Gain}]
\small
Untuk mengukur efektivitas pelatihan, hitung \textit{Normalized Learning Gain} menggunakan formula berikut:

\vspace{0.3cm}
\begin{center}
$\displaystyle g = \frac{\text{Skor Post-Test} - \text{Skor Pre-Test (Bag. A)}}{100 - \text{Skor Pre-Test (Bag. A)}} \times 100\%$
\end{center}

\vspace{0.3cm}
\renewcommand{\arraystretch}{1.3}
\begin{tabularx}{\textwidth}{@{} c X @{}}
\toprule
\textbf{Nilai $g$} & \textbf{Interpretasi} \\
\midrule
$g \geq 70\%$ & \textbf{Tinggi} -- Pelatihan sangat efektif; peserta mengalami perubahan signifikan \\
$30\% \leq g < 70\%$ & \textbf{Sedang} -- Pelatihan cukup efektif; masih ada ruang perbaikan metode \\
$g < 30\%$ & \textbf{Rendah} -- Evaluasi ulang metode penyampaian atau kesiapan prasyarat peserta \\
\bottomrule
\end{tabularx}
\renewcommand{\arraystretch}{1.0}

\vspace{0.3cm}
\textbf{Catatan:} Perbandingan hanya dilakukan pada soal Bagian A (pilihan ganda) yang memiliki padanan langsung dengan Post-Test. Bagian B dan C pre-test berfungsi sebagai data kualitatif untuk menyesuaikan pendekatan fasilitasi.
\end{plnbox}

\vspace{0.5cm}

\subsection{Peta Keterkaitan Pre-Test dengan Materi Modul}

\renewcommand{\arraystretch}{1.3}
\begin{center}
\small
\begin{tabularx}{0.95\textwidth}{@{} c X X @{}}
\toprule
\textbf{Soal Pre-Test} & \textbf{Prasyarat yang Diuji} & \textbf{Akan Diperdalam di Bab} \\
\midrule
A1 (PP 71/2019) & Regulasi PSTE & Bab 4 (SOP \& Regulasi), Bab 6 (Resistensi TTD Digital) \\
A2 (Definisi Arsip) & UU 43/2009 & Bab 1 (Konteks Strategis) \\
A3 (\textit{Bottleneck}) & Terminologi proses bisnis & Bab 2 (\textit{Friction Points} \& Evaluasi) \\
A4 (Metadata) & Konsep dasar metadata & Bab 2 (Diagnostik), Bab 4 (Metadata SOP) \\
A5 (JRA) & Retensi arsip & Bab 4 (Komponen Klasifikasi, Retensi \& Akses) \\
A6 (SOP) & Fungsi SOP & Bab 4 (Evolusi SOP Digital) \\
A7 (Integrasi) & Integrasi sistem & Bab 1 (Makna Integrasi), Bab 3 (Pola Integrasi) \\
A8 (TTD Elektronik) & Hukum TTD elektronik & Bab 4 (Regulasi), Bab 6 (\textit{Change Management}) \\
A9 (Siklus Hidup) & \textit{Records lifecycle} & Bab 2 (\textit{Value Stream Mapping}) \\
A10 (Risiko Audit) & Dampak tata kelola buruk & Bab 1 (Konteks Strategis), Bab 5 (KPI \& ROI) \\
\midrule
B1--B5 & Tingkat familiaritas topik & Menentukan kecepatan penyampaian per bab \\
\midrule
C1--C5 & Pengalaman nyata di unit kerja & Input studi kasus workshop Bab 2--4 \\
\bottomrule
\end{tabularx}
\end{center}
\renewcommand{\arraystretch}{1.0}

\end{document}
